\subsubsection{FFLabelClass}

The class {\bf FFLabelClass} is the superclass of all labels and
defines some methods that all labels must provide. For details on
what the methods should do in a label, see section
\vref{writing_labels}. The methods
{\bf FFLabelClass} defines are:

\begin{itemize}

\item{\bf void \_setType(const String\& s)}

Appends ":" followed by {\bf s} to the type. Should be called in
the constructor of each label.

\item{\bf virtual BOOL parse(const String\& str)}

Does nothing. Returns {\bf YES}.

\item{\bf virtual String unparse()}
  
Does nothing. Returns "".

\item{\bf virtual const String\& type() const}

Returns a reference to the type of the label.

\item{\bf virtual String baseType() const}

Returns the part after the last ":" in the type of the label.

\item{\bf virtual BOOL isSaveable() const}

Returns {\bf NO}. 

\item{\bf virtual BOOL shouldCopy() const}

Returns {\bf YES}.

\item{\bf virtual BOOL isInspectable() const}

Returns {\bf NO}.
  
\item{\bf virtual BOOL compare(const FFLabelClass*) const}

Returns {\bf NO}.

\item{\bf virtual void event(labelEvent, labelEventStruct *) }

Does nothing.

\item{\bf virtual void assign(const FFLabelClass*)}

Does nothing.

\end{itemize}

\subsubsection{FFGraphLabel}

The class {\bf FFGraphLabel} is a subclass of {\bf FFLabelClass}.
{\bf FFGraphLabel} is the superclass for all graph labels. It
defines the following methods:

\begin{itemize}
  
\item{\bf FFGraphLabel()}

Appends "FFGraphLabel" to the type.

\item{\bf virtual void init(const FFGraph\& g)}
  
  Tells the label that it is bound to graph {\bf g}. The member
  variable {\bf graph} is set to {\bf g}.

\item{\bf virtual FFGraphLabel *proto() const = NULL}

This method must be implemented by subclasses to reproduce itself.

\item{\bf virtual BOOL hasOwnInspector() const}

Returns {\bf YES} if the member variable {\bf inspector != NULL}.
 
\item{\bf virtual void showInspector()}
  
  If the graph label has an inspector, the methods {\bf
    FFInspector::show()} and {\bf FFInspector::inspect()} are
  invoked.

\item{\bf virtual void updateInspector()}
  
  If the graph label has an inspector, the method {\bf
    FFInspector::inspect()} is invoked.

\item{\bf virtual FFGraph getGraph()}
  
  Returns the member variable {\bf graph} which usually contains
  the graph, the label is bound to.

\item{\bf virtual \~FFGraphLabel()}

If the label has an inspector, the inspector is freed.

\end{itemize}

\subsubsection{FFNodeLabel}

The class {\bf FFNodeLabel} is a subclass of {\bf FFLabelClass}. 
{\bf FFNodeLabel} is the superclass for all node labels. It defines
the following methods:

\begin{itemize}

\item{\bf FFNodeLabel()}

Appends ":" followed by "FFNodeLabel" to the type. 

\item{\bf virtual void init( FFNodeClass *n)}

Sets the member variable {\bf node} to {\bf n}. This method
is used by the library to tell a label to which node it is
bound.

\item{\bf virtual FFNodeLabel *proto() const = NULL}

Must be implemented by subclasses to reproduce themselves.

\item{\bf virtual \~FFNodeLabel()}

Does nothing.

\end{itemize}

\subsubsection{FFEdgeLabel}

The class {\bf FFEdgeLabel} is a subclass of {\bf FFLabelClass}. 
{\bf FFEdgeLabel} is the superclass for all edge labels. It defines
the following methods:

\begin{itemize}

\item{\bf FFEdgeLabel()}

Appends ":" followed by "FFEdgeLabel" to the type of the label.

\item{\bf virtual void init( FFEdgeClass *e)}

Sets the member variable {\bf edge} to {\bf e}. This method is used
by the library to tell it to which edge it is bound.

\item{\bf virtual FFEdgeLabel *proto() const = NULL}

Must be implemented by subclasses to reproduce themselves.

\item{\bf virtual \~FFEdgeLabel()}

Does nothing.

\end{itemize}

\subsubsection{FFSubGraphLabel}

The class {\bf FFSubGraphLabel} is a subclass of {\bf
  FFNodeLabel}. {\bf FFSubGraphLabel} implements nested graphs.
This means a {\bf FFSubGraphLabel} contains a graph itself.
The classes {\bf FFTclWindow} and {\bf FFTclSubWindow} support
the display of nested graphs.

\begin{itemize}

\item{\bf virtual BOOL parse(const String\& s)}

{\bf s} should contain the file name of a sgraph file. The
graph from the file is load into the graph the label contains.

\item{\bf virtual String unparse()}

Returns the filename of the graph which the label contains.

\item{\bf virtual BOOL isInspectable() const}

Returns YES.
    
\item{\bf FFSubGraphLabel()}

Adds ":" followed by "FFSubGraphLabel" to the type of the label.

\item{\bf virtual BOOL isSaveable() const}

Returns YES.

\end{itemize}

\subsubsection{FFLayoutLabel}

The class {\bf FFLayout} is a subclass of {\bf FFGraphLabel} and
is the superclass for all labels that do any layout in the
distribution. The class does not actually do anything, but is
designed to define a common interface for all layout labels.
See section \vref{writing_layouters} for more information. The class
defines the following methods:

\begin{itemize}

\item{\bf FFLayoutLabel()}

Adds ":" followed by "FFLayoutLabel" to the type of the label.

\item{\bf virtual void layout(AVLSet \verb|<|FFNodeClass
    *\verb|>|\&)}

Does nothing.

\item{\bf virtual void layout()}

Does nothing.

\end{itemize}

\subsubsection{FFSugi}
\label{ffsugi}

The class {\bf FFSugi} is a subclass of {\bf FFLayoutLabel} and
implements a Sugiyama layout algorithm. The class defines the
following methods:

\begin{itemize}

\item{\bf FFSugi ()}

Creates a new instance of the {\bf FFSugi} class.

\item{\bf virtual void init(const FFGraph\& g)}

Used internally by the library to initialize the label.   

\item{\bf virtual BOOL isInspectable() const}

Returns YES.

\item{\bf virtual FFGraphLabel* proto() const}

Returns a new new instance of {\bf FFSugi}.
  
\item{\bf BOOL getUseStartNode ()}

Returns whether start-node information is used for the layout.

\item{\bf void setUseStartNode (BOOL flag)}

Sets whether start-node information should be used for the layout.

\item{\bf BOOL getKeepEdgePoints ()}

Returns whether edge points are kept in a local layout.

\item{\bf void setKeepEdgePoints (BOOL)}

Sets whether edge points should be kept during a local layout.
 
\item{\bf BOOL getUseGrid ()}

Returns whether a grid should be used to position nodes and edges.

\item{\bf void setUseGrid (BOOL)}

Sets whether a grid should be used to position nodes and edges.
 
\item{\bf BOOL getLeavesDown ()}

Returns whether leaves should be put on the lowest level.

\item{\bf void setLeavesDown (BOOL)}
  
Sets whether leaves should be put on the lowest level.

\item{\bf int getCrossingAlg ()}

Returns what algorithm is used to reduce crossings.

\item{\bf BOOL setCrossingAlg (int alg)}
  
  Specifies that algorithm {\bf alg} should be used to reduce
  crossings.

\item{\bf int getIteration1 ()}
  
  Returns the number of iterations for the reduce crossings
  algorithm.

\item{\bf BOOL setIteration1 (int n)}
  
  Sets the number of iterations for the reduce crossings
  algorithm to {\bf n}.

\item{\bf FFCoord getxoff ()}
  
  Returns the lowest x coordinate the graph will have after a
  layout.

\item{\bf BOOL setxoff (FFCoord x)}
  
  Sets the lowest x-coordinate the graph may have after a layout
  to {\bf x}.

\item{\bf FFCoord getyoff ()}

  Returns the lowest y coordinate the graph will have after a
  layout.

\item{\bf BOOL setyoff (FFCoord y)}

  Sets the lowest y-coordinate the graph may have after a layout
  to {\bf y}.

\item{\bf FFCoord getzoff ()}

  Returns the lowest z coordinate the graph will have after a
  layout.

\item{\bf BOOL setzoff (FFCoord z)}
  
  Sets the lowest z-coordinate the graph may have after a layout
  to {\bf z}.

\item{\bf FFCoord getxDistance ()}
  
  Returns the distance the nodes will have on the x-coordinate
  after a layout.

\item{\bf BOOL setxDistance (FFCoord x)}

  Sets the distance the nodes will have on the x-coordinate
  after a layout to {\bf x}.

\item{\bf FFCoord getyDistance ()}
  
  Returns the distance the nodes will have on the y-coordinate
  after a layout.

\item{\bf BOOL setyDistance (FFCoord)}

  Sets the distance the nodes will have on the y-coordinate
  after a layout to {\bf y}.

\item{\bf FFCoord getzDistance ()}
  
  Returns the distance the nodes will have on the z-coordinate
  after a layout.

\item{\bf BOOL setzDistance (FFCoord z)}

  Sets the distance the nodes will have on the z-coordinate
  after a layout to {\bf z}.
  
\item{\bf FFCoord getDummyOffset ()}
  
  Returns the percentage of the distance of edge points relative
  to the distance of nodes.

\item{\bf void setDummyOffset (FFCoord p)}
  
  Sets the percentage of the distance of edge points relative to
  the distance of nodes to {\bf p}.
  
\item{\bf int getDimension ()}

  Returns the number of dimensions the graph should be layouted in.
  May return 2 or 3.

\item{\bf void setDimension (int i)}
  
  Sets the number of dimensions the graph should be layouted to
  {\bf i}. {\bf i} must be 2 or 3.

\item{\bf void layout (AVLSet \verb|<|FFNodeClass *\verb|>|\& set)}
  
  Does a local layout on the subgraph defined by the nodes in
  {\bf set} and the edges between the nodes in the set.

\item{\bf void layout ()}
  
  Does a sugiyama layout on the graph.

\end{itemize}

\subsubsection{FFSpring}
\label{ffspring}

The class {\bf FFSpring} is a subclass of {\bf FFLayoutLabel} and
implements a spring-embedder layout algorithm. The class defines the
following methods:

\begin{itemize}

\item{\bf  FFSpring ()}

Creates a new instance of the {\bf FFSpring} class.

\item{\bf virtual void init(const FFGraph\& g)}

Used internally by the library to initialize the label.   

\item{\bf  virtual FFGraphLabel* proto() const}

Returns a new new instance of {\bf FFSpring}.
  
\item{\bf  virtual BOOL isInspectable() const}

Returns YES.

\item{\bf  BOOL    getAnimation ()}

Returns whether the layout should be animated.

\item{\bf  void    setAnimation (BOOL)}

Sets if the layout should be animated.

\item{\bf  BOOL    getUseEdgeWeight ()}

Returns whether the weight of edges should be used.

\item{\bf  void    setUseEdgeWeight (BOOL)}

Sets whether the weight of edges should be used.

\item{\bf  int     getMaxIteration ()}

Returns the maximum number of iterations.

\item{\bf  BOOL    setMaxIteration (int n)}

Sets the maximum number of iterations to {\bf n}.

\item{\bf  int     getAnimationIntervals ()}

Returns the number of iterations before displaying a result.

\item{\bf  BOOL    setAnimationIntervals (int n)}

Sets the number of iterations before displaying a result to {\bf n}.

\item{\bf  FFCoord getOptDistance ()}

Returns the optimal distance between two nodes.

\item{\bf  BOOL    setOptDistance (FFCoord x)}
  
  Sets the optimal distance between two nodes to {\bf x}. Returns
  {\bf YES}, if x is a valid distance.

\item{\bf  FFCoord getxoff ()}

  Returns the lowest possible x-coordinate.

\item{\bf  BOOL    setxoff (FFCoord x)}
  
  Sets the lowest possible x-coordinate to {\bf x}.

\item{\bf  FFCoord getyoff ()}

  Returns the lowest possible y-coordinate.

\item{\bf  BOOL    setyoff (FFCoord)}

  Sets the lowest possible y-coordinate to {\bf y}.

\item{\bf  FFCoord getzoff ()}

  Returns the lowest possible z-coordinate.

\item{\bf  BOOL    setzoff (FFCoord)}

  Sets the lowest possible z-coordinate to {\bf z}.

\item{\bf  BOOL    getDimension ()}
  
  Returns {\bf YES} if 3 dimensions are used and {\bf NO} if 2
  dimensions are used.

\item{\bf  void    setDimension (BOOL flag)}

  If {\bf flag == YES} 3 dimensions are used for layout. If {\bf flag == NO},
  2 dimensions are used.
  
\item{\bf  void    layout (AVLSet\verb|<|FFNodeClass *\verb|>|\& set)}
  
  Does a local layout for the subgraph defined by the nodes in
  {\bf set} and the edges between nodes in {\bf set}.

\item{\bf  void    layout ()}
  
  Does a layout of the whole graph.

\end{itemize}
